\documentclass[a4paper,11pt]{article}
\usepackage[T1]{fontenc}
\usepackage[utf8]{inputenc}
\usepackage[italian]{babel}
\usepackage{amsmath,amssymb}
\usepackage{graphicx}
\usepackage{geometry}
\usepackage{xcolor}
\usepackage{listings}
\geometry{margin=2.5cm}

\lstdefinelanguage{MiniZinc}{
  morekeywords={int,set,of,array,var,constraint,include,solve,satisfy,output,show},
  sensitive=true,
  morecomment=[l]{\%},
  morestring=[b]"
}

\lstset{
  language=MiniZinc,
  basicstyle=\ttfamily\small,
  keywordstyle=\color{blue}\bfseries,
  commentstyle=\color{teal!70!black}\itshape,
  stringstyle=\color{orange!70!black},
  numbers=left,
  numberstyle=\tiny\color{gray},
  stepnumber=1,
  numbersep=8pt,
  showstringspaces=false,
  frame=single,
  breaklines=true
}

\title{CSP: Cards 2}
\date{}

\begin{document}
\maketitle

\subsection*{1.1 Modellazione} 
Una formulazione CSP ad alto livello è la seguente.

\paragraph{Variabili}
Definiamo una variabile per ogni carta:
\[
X = \{x_{i} \mid i \in \{1, \ldots, N\}\},
\]
Dove $i$ rappresenta la posizione della carta nella soluzione e $x_{i}$ rappresenta il numero associato alla carta.

\[
P_1, P_2, P_3
\]
Rappresentano la posizione dei ``pivot'' in cui la sequenza si inverte (passa da crescente a decrescente o viceversa).

\paragraph{Domini}
\[
\forall x \in X:\quad D(x)=\{1,\ldots,M\}.
\]

\[ 
D(P_1) = \{2, \ldots, N-3\}
\]
\[
D(P_2) = \{3, \ldots, N-2\}
\]
\[
D(P_3) = \{4, \ldots, N-1\}
\]

\paragraph{Vincoli}

Vincolo di cardinalità (le occorrenze dei valori assegnati devono corrispondere a quelle della collezione iniziale):
\[
C_1 = \langle (x_1, \ldots, x_N), \left\{ (v_1, \ldots, v_N) \mathrel{\Big|} \forall k \in \{1, \ldots, M\}, |\{i \mid v_i = k\}| = |\{j \mid Cards_j = k\}| \right\} \rangle
\]

Vincolo di ordinamento dei pivot:
\[
C_2 = \langle (P_1, P_2, P_3), \{ (p_1, p_2, p_3) \mid p_1 < p_2 < p_3 \} \rangle
\]

Vincolo sulla distanza dei massimi locali:
\[
C_3 = \langle (P_1, P_3), \{ (p_1, p_3) \mid p_3 - p_1 = D \} \rangle
\]

Vincoli di monotonicità:
Per imporre l'andamento della sequenza in base alla posizione rispetto ai pivot, definiamo una famiglia di vincoli. Per ogni $i \in \{1, \ldots, N - 1\}$:
\[
C_{4,i} = \langle (x_i, x_{i+1}, P_1, P_2, P_3), \left\{ (v_i, v_{i+1}, p_1, p_2, p_3) \mathrel{\Bigg|} 
\begin{array}{l}
(i < p_1 \implies v_i < v_{i+1}) \land \\
(p_1 \le i < p_2 \implies v_i > v_{i+1}) \land \\
(p_2 \le i < p_3 \implies v_i < v_{i+1}) \land \\
(i \ge p_3 \implies v_i > v_{i+1})
\end{array}
\right\} \rangle
\]

L'insieme totale dei vincoli del CSP è:
\[
C = \{C_1, C_2, C_3\} \cup \{C_{4,i} \mid i \in \{1, \ldots, N-1\}\}
\]

\newpage

\subsection*{1.2 Istanziazione per (Cards, N, M, D) = ([1,1,2,2,3,3,4], 7, 4, 4)}

\paragraph{Variabili}
\[
X = (x_1, x_2, x_3, x_4, x_5, x_6, x_7)
\]
\[
P = (P_1, P_2, P_3)
\]

\paragraph{Domini}
\[
\forall i \in \{1, \ldots, 7\}:\quad D(x_i) = \{1, 2, 3, 4\}
\]
\[
D(P_1) = \{2, 3, 4\}, \quad D(P_2) = \{3, 4, 5\}, \quad D(P_3) = \{4, 5, 6\}
\]

\paragraph{Vincoli}

\[
C_1 = \langle (x_1, \ldots, x_7), \left\{
  \text{tutte le permutazioni di (1, 1, 2, 2, 3, 3, 4)}
\right\} \rangle
\]

\[
C_2 = \langle (P_1, P_2, P_3), \left\{ 
\begin{array}{l}
(2, 3, 4), (2, 3, 5), (2, 3, 6), \\
(2, 4, 5), (2, 4, 6), (2, 5, 6), \\
(3, 4, 5), (3, 4, 6), (3, 5, 6), \\
(4, 5, 6)
\end{array}
\right\} \rangle
\]

\[
C_3 = \langle (P_1, P_3), \{ (2, 6) \} \rangle
\]

Esplicitiamo il primo vincolo della famiglia $C_{4,i}$ per $i=1$. 
\[
C_{4,1} = \langle (x_1, x_2, P_1, P_2, P_3), \left\{ 
\begin{array}{l}
(1, 2, 2, 3, 4), (1, 2, 2, 3, 5), \dots, (1, 2, 4, 5, 6), \\
(1, 3, 2, 3, 4), (1, 3, 2, 3, 5), \dots, (1, 3, 4, 5, 6), \\
(1, 4, 2, 3, 4), (1, 4, 2, 3, 5), \dots, (1, 4, 4, 5, 6), \\
(2, 3, 2, 3, 4), (2, 3, 2, 3, 5), \dots, (2, 3, 4, 5, 6), \\
(2, 4, 2, 3, 4), (2, 4, 2, 3, 5), \dots, (2, 4, 4, 5, 6), \\
(3, 4, 2, 3, 4), (3, 4, 2, 3, 5), \dots, (3, 4, 4, 5, 6)
\end{array}
\right\} \rangle
\]
\newpage

\subsection*{1.3 Codifica in MiniZinc}
\begin{lstlisting}
include "alldifferent.mzn";
int: N;
int: M;
int: D;
array[1..N] of int: Cards;
array[1..N] of var 1..M: X;
var 2..N-3: P1;
var 3..N-2: P2;
var 4..N-1: P3;

constraint forall(v in 1..M)(
    sum(i in 1..N)(X[i] == v) == sum(j in 1..N)(Cards[j] == v)
);

constraint P1 < P2 /\ P2 < P3;

constraint P3 - P1 == D;

constraint forall(i in 1..N-1) (
    (i < P1 -> X[i] < X[i+1]) /\
    (i >= P1 /\ i < P2 -> X[i] > X[i+1]) /\
    (i >= P2 /\ i < P3 -> X[i] < X[i+1]) /\
    (i >= P3 -> X[i] > X[i+1])
);

solve satisfy;

output [show(X)];

\end{lstlisting}



\end{document}